\section*{Capítulo \ref{ch:b1:periodicity} --- Periodicidad}

\begin{solution}{exo:b1:periodicity:1}
\ce{Ca} $[\ce{Ar}]\,4\text{s}^2$: \omterm{def:b1:periodicity:table}{periodo} 4, \omterm{def:b1:periodicity:table}{grupo} 2, \omterm{def:b1:periodicity:table}{bloque} s. \ce{Fe}
$[\ce{Ar}]\,3\text{d}^6 4\text{s}^2$: \omterm{def:b1:periodicity:table}{periodo} 4, \omterm{def:b1:periodicity:table}{grupo} 8, \omterm{def:b1:periodicity:table}{bloque} d. \ce{Se}
$[\ce{Ar}]\,3\text{d}^{10} 4\text{s}^2 4\text{p}^4$: \omterm{def:b1:periodicity:table}{periodo} 4, \omterm{def:b1:periodicity:table}{grupo} 16,
\omterm{def:b1:periodicity:table}{bloque} p. \ce{I} $[\ce{Kr}]\,4\text{d}^{10} 5\text{s}^2 5\text{p}^5$:
\omterm{def:b1:periodicity:table}{periodo} 5, \omterm{def:b1:periodicity:table}{grupo} 17, \omterm{def:b1:periodicity:table}{bloque} p.
\end{solution}

\begin{solution}{exo:b1:periodicity:2}
\ce{Na} es más grande que \ce{Cl}: mismo \omterm{def:b1:periodicity:table}{periodo}, $Z^{*}$ menor. \ce{Na}
es más grande que \ce{Na+}, que ha perdido su capa 3s. \ce{Cl-}
(\qty{181}{pm}) es más grande que \ce{F-} (\qty{133}{pm}): una capa más.
\ce{Na+} (\qty{102}{pm}) es más grande que \ce{Mg^{2+}} (\qty{72}{pm}):
los mismos diez electrones, carga nuclear 11 frente a 12.
% ledger: rion:Cl-VI, rion:F-VI, rion:Na+VI, rion:Mg2+VI
\end{solution}

\begin{solution}{exo:b1:periodicity:3}
$\ce{Na} (5.14) < \ce{Al} (5.99) < \ce{Mg} (7.65) < \ce{S} (10.36) <
\ce{P} (10.49) < \ce{Ar} (15.76)$. El aumento general a lo largo del
\omterm{def:b1:periodicity:table}{periodo} se rompe dos veces: \ce{Al} por debajo de \ce{Mg} (el electrón
arrancado es el primer electrón 3p, de energía más alta que 3s) y \ce{S}
por debajo de \ce{P} (el electrón arrancado es el primer electrón 3p
apareado).
\end{solution}

\begin{solution}{exo:b1:periodicity:4}
La \omterm{def:b1:periodicity:polarisability}{polarizabilidad} mide la facilidad con que un campo eléctrico deforma
la nube electrónica. \ce{I-} es más polarizable que \ce{F-}: es más
grande y sus electrones externos están más lejos del núcleo. \ce{Cl-} (18
electrones retenidos por $Z = 17$) es mucho más polarizable que \ce{Na+}
(10 electrones retenidos por $Z = 11$), como lo son los aniones frente a
los cationes.
\end{solution}

\begin{solution}{exo:b1:periodicity:5}
Razones: $18.83/5.99 = 3.1$, $28.45/18.83 = 1.5$, $119.99/28.45 = 4.2$.
El salto llega después del tercer electrón: tres \omterm{def:b1:quantum-numbers:configuration}{electrones de valencia},
\omterm{def:b1:periodicity:table}{grupo} 13, \omterm{def:b1:periodicity:table}{periodo} 3: el aluminio. Forma \ce{Al^{3+}}, $1\text{s}^2 2\text{s}^2
2\text{p}^6$.
% ledger: ie:Al, ie2:Al, ie3:Al, ie4:Al
\end{solution}

\begin{solution}{exo:b1:periodicity:6}
El electrón añadido de \ce{F-} entra en la pequeña capa 2p, donde lo
repelen con fuerza los siete \omterm{def:b1:quantum-numbers:configuration}{electrones de valencia} que ya están en ella;
en la capa 3p, más grande, del cloro la repulsión es menor, de modo que se
libera más energía. En el nitrógeno ($2\text{p}^3$, \omorbs{u,u,u}) el
electrón adicional tendría que aparearse en una \omterm{def:b1:quantum-numbers:orbital}{subcapa} semillena; en el
magnesio ($3\text{s}^2$) tendría que inaugurar la \omterm{def:b1:quantum-numbers:orbital}{subcapa} 3p, más alta. En
ninguno de los dos casos el anión es más estable que el átomo y un
electrón libre.
\end{solution}

\begin{solution}{exo:b1:periodicity:7}
$\chi_M(\ce{Na}) = (5.14 + 0.55)/2 = \qty{2.85}{eV}$; $\chi_M(\ce{Cl}) =
(12.97 + 3.61)/2 = \qty{8.29}{eV}$. La diferencia es grande: el par de
electrones de un enlace \ce{Na-Cl} pertenece casi por completo al cloro, y
el cloruro de sodio es iónico, \ce{Na+ Cl-}.
\end{solution}

\begin{solution}{exo:b1:periodicity:8}
\ce{H-F}: $\Delta\chi = 1.78$, en el límite de la región iónica (en la
práctica, un enlace covalente muy polar), \ce{F^{\delta-}}. \ce{C-H}: 0.35,
casi apolar. \ce{Na-Cl}: 2.23, iónico, \ce{Cl^{-}}. \ce{C-Cl}: 0.61,
covalente polar, \ce{C^{\delta+}-Cl^{\delta-}}. \ce{O-H}: 1.24, covalente
polar, \ce{O^{\delta-}-H^{\delta+}}.
\end{solution}

\begin{solution}{exo:b1:periodicity:9}
Al bajar en el grupo, el electrón de valencia está en una capa de $n$
mayor, más lejos del núcleo y apantallado por más \omterm{def:b1:quantum-numbers:configuration}{electrones internos}; se
arranca con más facilidad. El rubidio, con una capa más, tiene una
\omterm{def:b1:periodicity:ionisation-energy}{primera energía de ionización} inferior a \qty{4.34}{eV}.
\end{solution}

\begin{solution}{exo:b1:periodicity:10}
El cinc es $[\ce{Ar}]\,3\text{d}^{10} 4\text{s}^2$: su electrón procede
de una \omterm{def:b1:quantum-numbers:orbital}{subcapa} 4s llena; el galio, $\dots 4\text{s}^2 4\text{p}^1$, pierde
su único electrón 4p, de energía más alta y apantallado por las \omterm{def:b1:quantum-numbers:orbital}{subcapas}
4s y 3d llenas. El arsénico, $4\text{p}^3$, pierde un electrón de una
\omterm{def:b1:quantum-numbers:orbital}{subcapa} semillena; el selenio, $4\text{p}^4$, pierde el electrón
apareado, desestabilizado por la repulsión de su compañero: su energía de
ionización es ligeramente menor.
\end{solution}

\begin{solution}{exo:b1:periodicity:11}
$\Delta E = (3.98 - 2.20)^2 = 1.78^2 = \qty{3.17}{eV} = 3.17 \times 96.5 =
\qty{306}{kJ/mol}$. Entonces, $D(\ce{H-F}) = \tfrac12(436 + 158) + 306 = 297 +
306 \approx \qty{603}{kJ/mol}$. El gran exceso es la atracción adicional
entre las cargas parciales $\ce{H^{\delta+}}$ y $\ce{F^{\delta-}}$, debida
a la gran diferencia de \omterm{def:b1:periodicity:electronegativity}{electronegatividad}.
\end{solution}

\begin{solution}{exo:b1:periodicity:12}
La \omterm{def:b1:periodicity:electronegativity}{escala de Pauling} se define a partir de las energías de los enlaces
A--B: un elemento que no forma ningún enlace no tiene valor de Pauling.
El helio, el neón y el argón no forman compuestos estables. El criptón y,
sobre todo, el xenón tienen capas de valencia más grandes y polarizables y
energías de ionización más bajas (\qty{14.0}{eV} para el criptón frente a
\qty{15.76}{eV} para el argón); el flúor y el oxígeno los oxidan y forman
compuestos como \ce{XeF2}, de modo que pueden medirse sus energías de
enlace.
% ledger: ie:Kr, ie:Ar
\end{solution}

\begin{solution}{pb:b1:periodicity:1}
\textbf{1.} Cada electrón se arranca de un ion más positivo y más pequeño
que el anterior, que retiene sus electrones con más fuerza.
\textbf{2.} $15.04/7.65 = 1.97$; $80.14/15.04 = 5.33$; $109.27/80.14 = 1.36$.
\textbf{3.} El salto llega después de dos electrones: X tiene dos
\omterm{def:b1:quantum-numbers:configuration}{electrones de valencia}, \omterm{def:b1:periodicity:table}{grupo} 2; en el \omterm{def:b1:periodicity:table}{periodo} 3 es el magnesio.
\textbf{4.} \ce{Mg^{2+}}, $1\text{s}^2 2\text{s}^2 2\text{p}^6$, la
configuración del neón.
\textbf{5.} Un metal: \omterm{def:b1:periodicity:table}{grupo} 2, a la izquierda de la tabla, con \omterm{def:b1:periodicity:ionisation-energy}{primeras
energías de ionización} bajas.
\textbf{6.} El tercer electrón procede del núcleo 2p, de $n$ menor, mucho
más cerca del núcleo y apenas apantallado, y se arranca de un ion con dos
cargas.
% ledger: ie:Mg, ie2:Mg, ie3:Mg, ie4:Mg
\textbf{7.} Todos tienen diez electrones: $1\text{s}^2 2\text{s}^2 2\text{p}^6$.
\textbf{8.} Los mismos diez electrones están retenidos por cargas
nucleares 8, 9, 11, 12 y 13: la nube se contrae a medida que crece $Z$.
\textbf{9.} \ce{O^{2-}}: el más grande, con la carga nuclear más baja
para sus diez electrones, y además un anión.
\textbf{10.} $140/54 = 2.6$.
\textbf{11.} $198.8/2 = \qty{99.4}{pm}$; el anión es $181/99.4 = 1.8$
veces más grande que el \omterm{def:b1:periodicity:radius}{radio covalente}.
% ledger: re:Cl2, rion:Cl-VI
\textbf{12.} $\ce{Mg^{2+}} < \ce{Na+} < \ce{Cl-}$: los dos cationes son
isoelectrónicos (carga 12 frente a 11), y \ce{Cl-} tiene una tercera capa.
\textbf{13.} $\chi_M$ = 10.41 (\ce{F}), 8.29 (\ce{Cl}), 7.59 (\ce{Br}),
7.53 (\ce{O}), 6.26 (\ce{C}), en eV.
\textbf{14.} Mulliken: $\ce{F} > \ce{Cl} > \ce{Br} > \ce{O} > \ce{C}$;
Pauling: $\ce{F} > \ce{O} > \ce{Cl} > \ce{Br} > \ce{C}$. El oxígeno pasa
del cuarto al segundo lugar.
\textbf{15.} $3.98/10.41 = 0.382$, $3.16/8.29 = 0.381$, $2.96/7.59 =
0.390$: casi constante, en torno a 0.38; para los halógenos, las dos
escalas son proporcionales.
\textbf{16.} La \omterm{def:b1:periodicity:electron-affinity}{afinidad electrónica} del oxígeno es pequeña
(\qty{1.44}{eV}) porque el electrón añadido debe aparearse en la compacta
capa 2p; el valor de Mulliken del átomo libre infravalora la atracción que
ejerce el oxígeno en sus enlaces, que la \omterm{def:b1:periodicity:electronegativity}{escala de Pauling}, construida a
partir de enlaces, mide directamente.
\textbf{17.} El cloro en \ce{HCl}; el flúor en \ce{ClF}.
\textbf{18.} $\Delta\chi = 0.96$: covalente polar.
\textbf{19.} $D(\ce{H-H}) = 2 \times 218.0 = \qty{436.0}{kJ/mol}$.
\textbf{20.} $D(\ce{Cl-Cl}) = 2 \times 121.3 = \qty{242.6}{kJ/mol}$.
\textbf{21.} $D(\ce{H-Cl}) = 218.0 + 121.3 - (-92.3) = \qty{431.6}{kJ/mol}$.
\textbf{22.} $\Delta E = 431.6 - \tfrac12(436.0 + 242.6) = 431.6 - 339.3 =
\qty{92.3}{kJ/mol}$. En símbolos, $D(\ce{H-Cl}) = \Delta_fH(\ce{H}) +
\Delta_fH(\ce{Cl}) - \Delta_fH(\ce{HCl})$, mientras que la media de los
enlaces simétricos es $\Delta_fH(\ce{H}) + \Delta_fH(\ce{Cl})$, de modo que $\Delta E =
-\Delta_fH(\ce{HCl})$.
\textbf{23.} $92.3/96.485 = \qty{0.957}{eV}$.
\textbf{24.} $|\chi(\ce{Cl}) - \chi(\ce{H})| = \sqrt{0.957} =
\textbf{0.98}$, frente a $3.16 - 2.20 = 0.96$ en las tablas: la
aritmética de Pauling, hecha con datos modernos, da la diferencia
tabulada con un margen de 0.02.
\end{solution}
